\documentclass[12pt]{article}
\usepackage[canadien]{babel}
\usepackage[utf8]{inputenc}
\usepackage[T1]{fontenc}
\usepackage{fancyhdr}
\usepackage{graphicx}
\usepackage{enumerate}
\usepackage{amsmath}
\usepackage{amssymb}
\usepackage{subfigure}
\usepackage{amsmath}
\usepackage{amssymb}
\usepackage{bm}
\usepackage{url}
\usepackage[dvipsnames]{xcolor}
\usepackage{todonotes}
\usepackage{hyperref}

% If your paper is accepted, change the options for the package
% aistats2e as follows:
%
%\usepackage[accepted]{aistats2e}
%
% This option will print headings for the title of your paper and
% headings for the authors names, plus a copyright note at the end of
% the first column of the first page.
\setlength{\parindent}{0cm}
\addtolength{\oddsidemargin}{-2cm}
\addtolength{\evensidemargin}{-2cm}
\setlength{\textwidth}{17.78cm}
\addtolength{\topmargin}{-2.25cm}
\setlength{\textheight}{23.24cm}
\addtolength{\parskip}{5mm}
\pagestyle{fancy}

%************
%* COMMANDS *
%************

\input{../math_commands.tex}

\newif\ifexercise
\exercisetrue
%\exercisefalse

\newif\ifsolution
\solutiontrue
%\solutionfalse

\usepackage{booktabs}
\usepackage[ruled,vlined]{algorithm2e}

\usepackage{amsthm}
\theoremstyle{definition}
\newtheorem{exercise}{Question}%[chapter]
\newtheorem{answer}{Answer} % asterisk to remove ordering

\newcommand{\Exercise}[1]{
\ifexercise#1\fi
}

\definecolor{answer}{rgb}{0.00, 0.12, 0.60}
\newcommand{\Answer}[1]{
\ifsolution\color{answer}\begin{answer}#1\end{answer}\fi
}

\newif\ifexercise
\exercisetrue
%\exercisefalse

\newif\ifsolution
\solutiontrue
%\solutionfalse

\usepackage{enumitem}
\newcommand{\staritem}{
\addtocounter{enumi}{1}
\item[$\phantom{x}^{*}$\theenumi]}
\setlist[enumerate,1]{leftmargin=*, label=\arabic*.}

\begin{document}


\fancyhead{}
\fancyfoot{}

\fancyhead[L]{
  \begin{tabular}[b]{l}
    IFT6135-A2026  \\
    Prof.~: Aishwarya Agrawal \\
  \end{tabular}
}
\fancyhead[R]{
  \begin{tabular}[b]{r}
    Devoir 2 - Programmation  \\
    Réseaux de neurones récurrents et Transformers  \\
  \end{tabular}
}
\fancyfoot[C]{- Ne pas distribuer -}

\vspace{1cm}

\shorthandoff{:}
{\textbf{Date de remise~: 15 novembre 2026 à 23 h}}\\


\vspace{-0.5cm}
\underline{Instructions}%
\renewcommand{\labelitemi}{\textbullet}

\renewcommand{\labelitemi}{\textbullet}

\begin{itemize}
\item \emph{Ce devoir est exigeant, commencez bien à l'avance~!}
\item \emph{Pour toutes les questions, montrez votre démarche.}
% Link to hw2_student.ipynb (both questions). To keep this link, upload hw2_student.ipynb as a new version of this
% Drive file (right-click > Manage versions > Upload new version); otherwise paste the new link here.
\item \emph{La partie pratique est fournie sous la forme d'un seul notebook Jupyter sur Google Colab, qui contient les
deux questions~} \emph{. Enregistrez-en une copie dans votre propre Drive avant de le modifier, et exécutez-le sur un GPU
(Runtime $\rightarrow$ Change runtime type $\rightarrow$ T4 GPU). Remettez ce notebook (File $\rightarrow$ Download
$\rightarrow$ Download .ipynb) sur Gradescope: les questions 1 et 2 sont toutes deux
corrigées à partir de celui-ci.}
\item \emph{Vous devez aussi remettre votre rapport (PDF) sur Gradescope. Votre rapport doit contenir les réponses
aux questions 1 et 2. Vous devez également fournir votre notebook Google Colab, avec les sorties déjà générées et un
accès public.}
\item \emph{Les auxiliaires d'enseignement pour ce devoir sont \textbf{Qichen Huang} et \textbf{Serena Chan}.}
\end{itemize}
\vspace{0.2cm}


\vspace{0.2cm}

\textbf{Lien vers le notebook: }\href{https://colab.research.google.com/drive/1U3l83ZcKE-iTDYsCHJfDzjuP3okoZihf?usp=sharing}{Lien}

% Q1
\begin{exercise}[50] \textbf{Apprendre à mettre à jour une mémoire avec des LSTM}
\Exercise{\label{ex:lstm}

Dans cette question, vous allez implémenter un LSTM et étudier s'il peut retenir la dernière valeur d'une clé après
des mises à jour répétées et des entrées non pertinentes. Vous utiliserez PyTorch. La différentiation automatique,
\texttt{nn.Linear}, \texttt{nn.Embedding} et les opérations primitives sur les tenseurs sont permises. N'utilisez pas
\texttt{nn.LSTM}, \texttt{nn.LSTMCell}, d'autres modules récurrents intégrés, ni une recherche dans un dictionnaire à
l'intérieur de votre modèle.

\textbf{Fichiers fournis.} La question 1 est la première partie du notebook du devoir 2 (lien dans les
instructions), qui explique la démarche. Complétez uniquement les blocs TODO indiqués de \texttt{lstm.py}, dans la
cellule du notebook qui commence par \texttt{\%\%writefile lstm.py}. Ne modifiez ni ne supprimez cette première
ligne~: le correcteur automatique lit votre code dans cette cellule. \texttt{memory\_data.py}, \texttt{train.py} et
\texttt{test\_public.py} sont fournis. Conservez les signatures des fonctions, les noms des paramètres et les formats
des tenseurs. Exécutez la cellule \texttt{lstm.py} pour créer le fichier, puis exécutez dans l'ordre les tests publics
et les cellules d'expérience.

\subsection*{Tâche et données}
Il y a trois clés $A$, $B$ et $C$, et quatre valeurs possibles~: $0$, $1$, $2$ et $3$.
Une mise à jour comme $A=2$ est représentée par \emph{un seul} jeton~; les clés et les valeurs ne sont pas des jetons
successifs distincts. Une séquence contient six jetons de mise à jour, suivis de $D$ jetons distracteurs et d'un
jeton de requête. Chaque distracteur est tiré parmi quatre symboles dédiés qui ne portent aucune mise à jour. La
cible est la valeur la plus récente attribuée à la clé interrogée. Par exemple,
\[
 A=1,\ B=0,\ C=2,\ A=3,\ B=1,\ C=0,\ D_0,\ D_2,\ Q_A
 \quad\longrightarrow\quad 3.
\]
L'ancienne valeur $1$ est incorrecte, tout comme $0$, la valeur de la dernière mise à jour d'une autre clé. La
requête n'est présentée qu'à la fin de la séquence.

Avec une clé $k$ et une valeur $v$ indexées à partir de zéro, les jetons de mise à jour ont l'identifiant $1+4k+v$~;
les jetons de requête ont l'identifiant $13+k$~; les identifiants des distracteurs sont $16$, $17$, $18$ et $19$.
L'identifiant $0$ est réservé au remplissage à droite, de sorte que la taille du vocabulaire est $20$. Les étiquettes
de valeur sont $0$, $1$, $2$ et $3$. Le modèle reçoit uniquement les identifiants des jetons et les longueurs valides
des séquences, jamais les métadonnées d'analyse du jeu de données.

Le générateur fourni crée 6\,000 exemples d'entraînement et 1\,200 exemples de validation. L'entraînement et la
validation utilisent $D\in\{0,2,4,8\}$ et $R\in\{1,2,3\}$ mises à jour de la clé interrogée. Sa dernière mise à jour
est à la position 2, 3 ou 4 (indexée à partir de zéro). Les autres mises à jour concernent d'autres clés~; en
particulier, le dernier jeton de mise à jour concerne toujours une clé différente. Les mises à jour consécutives de la
clé interrogée ont des valeurs différentes. Le générateur équilibre les paires (clé interrogée, valeur cible) et
rejette les séquences de jetons en double, au sein des ensembles fournis et entre eux. En entraînement et en
validation, les positions de la dernière mise à jour sont approximativement équilibrées dans l'ensemble, mais pas
forcément au sein de chaque condition $(D,R)$.

\subsection*{Partie 1~: Implémentation du modèle (30 points)}
Pour une taille de lot $B$, une dimension de plongement $E$ et une dimension cachée $H$, l'entrée de la cellule est
$x_t\in\mathbb{R}^{B\times E}$ et les deux états sont de forme $B\times H$. Utilisez l'ordre des portes
$(i,f,g,o)$~:
\[
 [a_i,a_f,a_g,a_o]=[x_t,h_{t-1}]W^\top+b,
 \qquad i_t=\sigma(a_i),\ f_t=\sigma(a_f),\ g_t=\tanh(a_g),\ o_t=\sigma(a_o),
\]
\[
 c_t=f_t\odot c_{t-1}+i_t\odot g_t,
 \qquad h_t=o_t\odot\tanh(c_t).
\]
\begin{enumerate}
\item \textbf{Cellule LSTM (18 points).} Complétez \texttt{LSTMCell.forward}~: calculez les entrées affines des
portes dans l'ordre spécifié (6), appliquez les non-linéarités appropriées (6) et retournez les états caché et de
cellule mis à jour (6). La couche affine est déjà instanciée. Votre implémentation doit prendre en charge les entrées
en lot et préserver l'autograd, le type de données (\emph{dtype}) de l'entrée et le périphérique (\emph{device}).

\item \textbf{Modèle de séquence (12 points).} Complétez \texttt{MemoryLSTM.recurrent} (8) et
\texttt{MemoryLSTM.forward} (4). Réinitialisez les états caché et de cellule à zéro à chaque appel de
\texttt{forward}. Traitez les pas de temps valides dans l'ordre. Aux positions de remplissage à droite, conservez
\emph{les deux} états précédents. Retournez l'état caché après le dernier jeton valide, requête comprise. Combinez le
plongement, la cellule récurrente et le classificateur fournis pour retourner quatre logits bruts par exemple, sans
softmax final. Le modèle utilise une seule couche récurrente, $E=32$ et $H=64$.
\end{enumerate}

Exécutez \texttt{python -m unittest test\_public -v} avant l'entraînement. Les tests publics sont illustratifs~; des
tests distincts vérifieront les valeurs numériques, les gradients, le traitement des lots et le comportement avec le
remplissage.

\subsection*{Partie 2~: Entraînement et analyse de la mémoire (20 points)}
Utilisez le script d'entraînement fourni et la graine 6135. La configuration est de 20 époques, une taille de lot de
128, Adam avec un taux d'apprentissage de 0,003 et un écrêtage de la norme du gradient à 1,0. N'ajoutez pas de couches
et ne faites pas de recherche d'hyperparamètres. Sélectionnez le point de contrôle (\emph{checkpoint}) qui a la plus
faible entropie croisée de validation, en gardant la première des époques en cas d'égalité. N'utilisez pas d'ensemble
de test pour sélectionner un point de contrôle. L'exécution sur CPU est possible. Le script d'évaluation fourni
utilise \texttt{model.eval()} et \texttt{torch.no\_grad()} sans mettre à jour les paramètres du modèle.

\begin{enumerate}
\item \textbf{Entraînement et évaluation (6 points).} Rapportez les courbes de perte et d'exactitude d'entraînement et
de validation (2), l'époque sélectionnée et la règle de sélection (2), ainsi que le matériel, les versions des
logiciels, le temps total d'entraînement et de validation et le nombre de paramètres du modèle (2). Utilisez
directement les sorties fournies par le notebook~; aucune expérience supplémentaire n'est requise. Conservez les
sorties de l'exécution.

\item \textbf{Analyse de la mémoire (8 points).} Évaluez le même point de contrôle sélectionné, sans réentraînement,
sur les 15 conditions fournies $D\in\{0,4,8,16,32\}$ et $R\in\{1,2,3\}$, avec 360 exemples par condition. Ici, la
dernière mise à jour de la clé interrogée est équilibrée entre les positions 2, 3 et 4 dans chaque condition. Sa
distribution de positions est donc la même lorsqu'on compare différentes valeurs de $R$, sans révéler une position de
réponse fixe. Rapportez un tableau d'exactitude et tracez l'exactitude en fonction de $D$, avec une courbe par valeur
de $R$ (3). Pour $R\geq2$, rapportez aussi le \emph{taux de valeur périmée} (\emph{stale-value rate})~: la proportion
de tous les exemples de cette condition dont la prédiction est égale à la valeur immédiatement précédente de la clé
interrogée (2). Indiquez que ce taux ne s'applique pas pour $R=1$. En un court paragraphe (environ 3 à 5 phrases),
expliquez les tendances et distinguez la performance observée d'une capacité de mémoire garantie (3). Indiquez
quelles longueurs de distracteurs étaient absentes de l'entraînement. Ne supposez pas que chaque augmentation de $D$
ou de $R$ doit réduire l'exactitude. Les différences entre les valeurs de $R$ doivent être interprétées de façon
descriptive, puisque les motifs et l'espacement des mises à jour antérieures varient aussi. Ces résultats portent sur
des distracteurs placés après toutes les mises à jour, et non sur des distracteurs intercalés entre les mises à jour.

\item \textbf{Analyse des échecs (6 points).} Montrez deux erreurs réelles lorsqu'il en existe au moins deux dans
l'ensemble des conditions de test, en utilisant les identifiants d'exemples exportés, les séquences décodées, les
vraies étiquettes, les prédictions et les valeurs précédentes (3). Privilégiez des types d'erreurs différents lorsque
c'est possible~; les deux types ne sont pas requis. Distinguez les erreurs de valeur périmée des autres erreurs et
donnez une explication plausible liée à la rétention et à la mise à jour de la mémoire, en 2 ou 3 phrases pour chaque
exemple (3). Vous n'avez pas à prouver la cause interne. S'il existe moins de deux erreurs, montrez toutes les erreurs
disponibles et indiquez leur nombre~; s'il n'y en a aucune, montrez deux exemples corrects et discutez des limites de
cette preuve. Ne faire aucune erreur ne fait pas perdre de points.
\end{enumerate}

\textbf{Remise.} Remettez le notebook (.ipynb) sur Gradescope~: la partie 1 est corrigée automatiquement à partir de
la cellule \texttt{\%\%writefile lstm.py}. Répondez à la partie 2 dans votre rapport PDF, en y incluant les
graphiques, le tableau des résultats et les exemples produits par le notebook.
Le travail expérimental est évalué selon la justesse, la reproductibilité et l'interprétation, et non selon l'atteinte
d'une exactitude prescrite.
}
\end{exercise}


% Q2
\begin{exercise}[50] \textbf{Apprentissage de représentations au niveau des mots et au niveau des phrases avec des Transformers}
\Exercise{\label{ex:transformer}}
\sloppy

\textbf{Tâche~:} Dans cette question, vous allez implémenter un encodeur Transformer bidirectionnel à partir de ses
équations avec PyTorch, puis en entraîner trois copies sur les mêmes phrases avec trois objectifs différents~: la
modélisation du langage masqué (MLM), la classification en inférence en langage naturel (NLI) et l'apprentissage
contrastif NLI. Vous comparerez ensuite ce que chaque objectif apprend au modèle, au niveau des mots individuels et
au niveau des phrases entières.

\textbf{Remise~:} La question 2 est la deuxième partie du notebook du devoir 2 (lien dans les instructions).
Complétez uniquement les blocs \texttt{\# TODO} dans les cellules dont la première ligne est
\texttt{\# GRADED CELL}~; tout le reste est fourni. Ne modifiez ni ne supprimez les lignes \texttt{\# GRADED CELL}, et
ne changez aucun nom de fonction, nom de paramètre ni format de tenseur. Sur Gradescope, remettez le notebook
(le même fichier que pour la question 1), qui est corrigé automatiquement pour les parties 1 et 2, et répondez à la
partie 3 dans votre rapport PDF.

\textbf{Données.}

Les trois modèles sont entraînés sur le corpus AllNLI (SNLI + MultiNLI)~: 1\,147\,270 phrases anglaises distinctes
d'au plus 64 jetons (12{,}6 en moyenne). Chaque phrase est mise en minuscules et découpée en mots, et chaque mot est
remplacé par son identifiant dans un vocabulaire des $V = 20\,000$ mots les plus fréquents. Les identifiants 0 à 3
sont réservés à \texttt{[PAD]}, \texttt{[MASK]}, \texttt{[CLS]} et \texttt{[UNK]} (mots hors vocabulaire), et chaque
phrase commence par \texttt{[CLS]}. Les trois objectifs utilisent les mêmes phrases sous des formes différentes~:
\begin{itemize}
  \item \textbf{MLM} utilise les phrases seules, sans étiquettes.
  \item \textbf{La classification NLI} utilise 942\,069 paires (prémisse, hypothèse, étiquette), où l'étiquette indique
        si l'hypothèse découle de la prémisse~: \emph{entailment} (implication), \emph{neutral} (neutre) ou
        \emph{contradiction}.
  \item \textbf{L'apprentissage contrastif NLI} utilise 557\,850 triplets (ancre, positif, négatif)~: une prémisse,
        une hypothèse qu'elle implique et une hypothèse qu'elle contredit.
\end{itemize}
Pour l'évaluation, nous utilisons 5\,000 phrases d'AllNLI mises de côté, sur lesquelles aucun modèle n'a été
entraîné, ainsi que le jeu de référence STS-B~: 1\,379 paires de phrases, chacune notée par des humains de 0 (sens
différent) à 5 (même sens).

\textbf{Notation.}

Un lot contient $B$ phrases complétées par du remplissage (\emph{padding}) jusqu'à une même longueur $T$.
$x \in \{0, \dots, V-1\}^{B \times T}$ sont les identifiants des jetons et $a \in \{0, 1\}^{B \times T}$ le masque
d'attention ($a_{b,t} = 1$ aux vrais jetons, $0$ aux \texttt{[PAD]}). L'encodeur a une dimension cachée $D = 128$,
$H = 4$ têtes d'attention de dimension $d = D/H = 32$ et $L = 4$ blocs. Il associe à chaque phrase un vecteur caché
par jeton, $O = \mathrm{Enc}(x) \in \mathbb{R}^{T \times D}$, de lignes $o_1, \dots, o_T$. Les trois objectifs
utilisent cette sortie différemment~:
\begin{itemize}
  \item \textbf{MLM} cache certains mots et prédit chacun d'eux à partir de son contexte, à partir du vecteur caché
        à sa position (une prédiction par mot).
  \item \textbf{La classification NLI} fait la moyenne des vecteurs cachés de chaque phrase pour obtenir un seul
        plongement de phrase (\emph{mean pooling}),
        \[
          z = \frac{\sum_{t=1}^{T} a_t\, o_t}{\sum_{t=1}^{T} a_t} \in \mathbb{R}^{D},
        \]
        puis prédit l'étiquette NLI d'une paire à partir des deux plongements, au moyen d'un petit classificateur
        (une prédiction par paire de phrases).
  \item \textbf{L'apprentissage contrastif NLI} utilise les mêmes plongements de phrases et entraîne chaque phrase à
        être plus proche, en similarité cosinus, de la phrase qu'elle implique que de toutes les autres phrases du
        lot (une prédiction par phrase).
\end{itemize}
La similarité cosinus de deux vecteurs $u, v \in \mathbb{R}^D$ est
$\cos(u, v) = \dfrac{u^\top v}{\|u\|\,\|v\|} \in [-1, 1]$~: elle vaut 1 quand $u$ et $v$ pointent dans la même
direction, 0 quand ils sont perpendiculaires (c.-à-d.\ sans rapport) et $-1$ quand ils pointent dans des directions opposées.

\textbf{Partie 1~: Implémentation du modèle (30 pts).}

\begin{enumerate}
  \item \textbf{Mise en lots et remplissage (2 pts).} Implémentez \texttt{collate(seqs, device)}. Étant donné $B$
  séquences d'identifiants $s_1, \dots, s_B$ de longueurs $n_1, \dots, n_B$, avec $T = \max_b n_b$, retournez
  \[
    x_{b,t} = \begin{cases} s_{b,t} & t \le n_b \\ 0 \ (\texttt{[PAD]}) & t > n_b \end{cases}, \qquad
    a_{b,t} = \mathbf{1}[t \le n_b].
  \]

  \item \textbf{Plongements (4 pts).} Implémentez \texttt{sinusoidal\_positions(T, D)}, qui retourne les encodages
  positionnels fixes $PE \in \mathbb{R}^{T \times D}$~:
  \[
    PE_{pos,\,2i} = \sin\!\left(\frac{pos}{10000^{2i/D}}\right), \qquad
    PE_{pos,\,2i+1} = \cos\!\left(\frac{pos}{10000^{2i/D}}\right),
  \]
  pour les positions $pos = 0, \dots, T-1$ et $i = 0, \dots, \tfrac{D}{2}-1$.
  Implémentez ensuite \texttt{Embeddings.forward(ids)}, l'entrée du premier bloc. Avec
  $E \in \mathbb{R}^{V \times D}$ la matrice des plongements de jetons et $E_{x_t}$ sa ligne pour le jeton $x_t$,
  \[
    X_t = \mathrm{Dropout}\big(\sqrt{D}\, E_{x_t} + PE_t\big), \qquad X \in \mathbb{R}^{T \times D}.
  \]

  \item \textbf{Attention multi-têtes et masque bidirectionnel (7 pts).} Implémentez
  \texttt{MultiHeadAttention.forward(x, mask)}. Chaque tête $h = 1, \dots, H$ a ses propres projections de requête,
  de clé et de valeur, qui vous sont fournies sous forme de listes de couches \texttt{nn.Linear(D, d)}
  \texttt{W\_q[h]}, \texttt{W\_k[h]}, \texttt{W\_v[h]}, et la projection de sortie \texttt{W\_o} est une
  \texttt{nn.Linear(D, D)}~:
  \[
    Q_h = W_Q^{(h)}(X), \quad K_h = W_K^{(h)}(X), \quad V_h = W_V^{(h)}(X) \in \mathbb{R}^{T \times d},
  \]
  \[
    A_h = \mathrm{Dropout}\!\left(\mathrm{softmax}\!\left(\frac{Q_h K_h^\top}{\sqrt{d}} + M\right)\right) V_h
    \in \mathbb{R}^{T \times d},
  \]
  \[
    \mathrm{MHA}(X) = W_O\big([A_1, \dots, A_H]\big) \in \mathbb{R}^{T \times D},
  \]
  où le softmax est calculé sur les clés (chaque ligne somme à 1), $[A_1, \dots, A_H]$ concatène les têtes selon la
  dimension des caractéristiques, et $M \in \mathbb{R}^{T \times T}$ est le masque d'attention~: $M_{t,s} = 0$ si la
  requête $t$ peut porter attention à la clé $s$, et $M_{t,s} = -10^4$ sinon (son poids après le softmax devient
  $\approx 0$).


  Notre encodeur est bidirectionnel~: chaque jeton peut porter attention à tous les autres jetons, avant ou après
  lui. La seule restriction est le remplissage, auquel on ne doit jamais porter attention. Implémentez
  \texttt{bidirectional\_mask(attention\_mask)}, qui retourne le masque booléen
  $m \in \{0, 1\}^{B \times T \times T}$ avec $m_{b,t,s} = a_{b,s}$ ($M_{t,s} = 0$ là où $m_{b,t,s} = 1$).

  \item \textbf{Bloc d'encodeur et encodeur (8 pts).} Implémentez \texttt{EncoderBlock.forward(x, mask)}, un bloc
  Transformer à post-normalisation~:
  \[
    X \leftarrow \mathrm{LN}_1\big(X + \mathrm{Dropout}(\mathrm{MHA}(X))\big), \qquad
    X \leftarrow \mathrm{LN}_2\big(X + \mathrm{Dropout}(\mathrm{FFN}(X))\big),
  \]
  où la sous-couche à propagation avant $\mathrm{FFN}(X) = W_2\, \mathrm{GELU}(W_1 X)$, de dimension cachée 512, est
  fournie.


  Implémentez ensuite \texttt{Encoder.forward(ids, attention\_mask, return\_logits)}~: construisez le masque,
  calculez les plongements des identifiants, appliquez les $L$ blocs dans l'ordre pour obtenir $O$, puis la tête MLM.
  Ainsi, le score du mot $w$ à la position $t$ est
  \[
    \ell_{t,w} = o_t^\top E_w + c_w, \qquad \ell = O E^\top + c \in \mathbb{R}^{T \times V},
  \]
  où $c \in \mathbb{R}^V$ est un biais appris. Un mot obtient un score élevé quand le vecteur caché est proche du
  plongement de ce mot.

  \item \textbf{Modélisation du langage masqué (9 pts).} Implémentez
  \texttt{mask\_tokens(ids, attention\_mask, vocab\_size, p)} comme suit~:
  \begin{enumerate}
    \item Sélectionnez chaque vrai jeton indépendamment avec probabilité $p = 0{,}15$, jamais un \texttt{[PAD]} ni
          un jeton spécial (identifiants 0 à 3, p.\ ex.\ \texttt{[CLS]}). Soit $\mathcal{M}$ l'ensemble des
          positions $(b, t)$ sélectionnées.
    \item Pour chaque position sélectionnée, indépendamment~: avec probabilité $0{,}8$, remplacez-la par
          \texttt{[MASK]}~; avec probabilité $0{,}1$, remplacez-la par un identifiant de mot tiré uniformément dans
          $\{4, \dots, V-1\}$~; avec probabilité $0{,}1$, laissez-la inchangée.
    \item Retournez les identifiants corrompus $\tilde{x}$ et les étiquettes $y$, où $y_{b,t} = x_{b,t}$ aux
          positions sélectionnées et $y_{b,t} = -100$ (ignorée) ailleurs.
  \end{enumerate}


  Implémentez ensuite \texttt{cross\_entropy(logits, targets, ignore\_index)} à la main (les trois objectifs
  l'utilisent). Pour des scores $z \in \mathbb{R}^{N \times C}$ et des cibles $y \in \{0, \dots, C-1\}^N$, en
  moyenne sur les lignes $\mathcal{K} = \{n : y_n \ne -100\}$~:
  \[
    \mathcal{L}_{\text{CE}}(z, y) = -\frac{1}{|\mathcal{K}|} \sum_{n \in \mathcal{K}}
      \left( z_{n, y_n} - \log \sum_{c=1}^{C} e^{z_{n,c}} \right).
  \]
  Utilisez \texttt{torch.logsumexp} pour le log de la somme. Implémentez ensuite
  \texttt{mlm\_loss(logits, labels)}, l'entropie croisée sur les positions masquées seulement. Avec les logits
  $\ell \in \mathbb{R}^{B \times T \times V}$ calculés à partir des identifiants corrompus $\tilde{x}$~:
  \[
    \mathcal{L}_{\text{MLM}} = -\frac{1}{|\mathcal{M}|} \sum_{(b,t) \in \mathcal{M}}
      \log \mathrm{softmax}(\ell_{b,t})_{\,y_{b,t}}.
  \]

  \item \textbf{Objectifs au niveau des phrases (fournis).} Lisez et comprenez les fonctions fournies, que vous utiliserez
  dans la partie 2~:
  \begin{itemize}
    \item \texttt{mean\_pool} et \texttt{sentence\_embedding}~: le plongement de phrase $z$ défini plus haut.
    \item \texttt{pair\_features(u, v)}~: l'entrée du classificateur NLI pour un plongement de prémisse $u$ et un
          plongement d'hypothèse $v$.
    \item \texttt{info\_nce(z\_a, z\_pos, z\_neg, tau)}~: la perte InfoNCE avec négatifs du lot et négatifs
          difficiles. Pour des ancres $z^a_i$, des positifs $z^+_i$ (implications) et des négatifs difficiles
          $z^-_i$ (contradictions), $i = 1, \dots, B$, et une température $\tau = 0{,}05$~:
          \[
            \mathcal{L}_{\text{con}} = -\frac{1}{B} \sum_{i=1}^{B} \log
              \frac{\exp\!\big(\cos(z^a_i, z^+_i)/\tau\big)}
                   {\sum_{j=1}^{B} \exp\!\big(\cos(z^a_i, z^+_j)/\tau\big)
                    + \sum_{j=1}^{B} \exp\!\big(\cos(z^a_i, z^-_j)/\tau\big)}.
          \]
  \end{itemize}
\end{enumerate}

\textbf{Partie 2~: Pertes et entraînement (6 pts).}

La boucle d'entraînement (optimiseur, journalisation, validation) est fournie. Implémentez~:

\begin{enumerate}
  \item \textbf{Étape d'entraînement (4 pts).} Implémentez \texttt{compute\_loss(model, batch, task, vocab\_size, tau)},
  qui retourne la perte d'un lot~:
  \begin{itemize}
    \item \texttt{task = "mlm"}~: masquez le lot avec \texttt{mask\_tokens}, encodez les identifiants
          corrompus et retournez $\mathcal{L}_{\text{MLM}}$.
    \item \texttt{task = "nli\_cls"}~: calculez les plongements de phrases $u$ des prémisses et $v$ des hypothèses,
          et retournez $\mathcal{L}_{\text{CE}}\big(\texttt{nli\_head}([u;\, v;\, |u - v|]),\ y\big)$, avec les
          étiquettes $y \in \{0, 1, 2\}$ (implication, neutre, contradiction).
    \item \texttt{task = "contrastive"}~: calculez les plongements de phrases des ancres, des positifs et des
          négatifs, et retournez $\mathcal{L}_{\text{con}}$.
  \end{itemize}

  \item \textbf{Calendrier du taux d'apprentissage (2 pts).} Implémentez
  \texttt{lr\_schedule(step, total\_steps, warmup\_steps)}, un préchauffage linéaire suivi d'une décroissance en
  cosinus. À l'étape $s = 0, 1, \dots, S - 1$, avec $S$ étapes au total et $W$ étapes de préchauffage, le taux
  d'apprentissage est $\eta \cdot m(s)$, où $\eta$ est le taux d'apprentissage maximal et
  \[
    m(s) = \begin{cases}
      \dfrac{s + 1}{W} & \text{si } s < W, \\[2ex]
      \dfrac{1}{2}\left(1 + \cos\!\left(\pi\, \dfrac{s - W}{S - W}\right)\right) & \text{sinon.}
    \end{cases}
  \]
\end{enumerate}

Entraînez ensuite les trois modèles avec la configuration donnée dans le notebook (environ 5 minutes chacun sur un
T4)~: chacun part d'une initialisation aléatoire et utilise AdamW (taux d'apprentissage maximal
$\eta = 5 \times 10^{-4}$, décroissance des poids (\emph{weight decay}) $0{,}01$, taille de lot 128, $W = 5\,\%$ des
étapes, écrêtage du gradient à une norme de 1{,}0).
\begin{itemize}
  \item \textbf{MLM~:} entraîné avec $\mathcal{L}_{\text{MLM}}$ pendant 6\,000 étapes.
  \item \textbf{Classification NLI~:} entraîné avec $\mathcal{L}_{\text{CE}}$ sur les étiquettes NLI pendant
        5\,000 étapes.
  \item \textbf{Contrastif NLI~:} entraîné avec $\mathcal{L}_{\text{con}}$ ($\tau = 0{,}05$) pendant 4\,000 étapes.
\end{itemize}
Un encodeur non entraîné de même architecture (initialisation aléatoire) est conservé comme point de comparaison.

\textbf{Partie 3~: Évaluation (14 pts).} Répondez dans votre rapport PDF. Toutes les métriques, tous les tableaux et
tous les graphiques sont calculés par des fonctions fournies dans le notebook. Incluez les sorties demandées par
chaque question.

Le notebook compare les quatre modèles selon deux métriques~:
\begin{itemize}
  \item \textbf{Exactitude sur les mots masqués}~: 15\,\% des mots des 5\,000 phrases mises de côté sont masqués
        comme dans la partie 1, et l'exactitude est la proportion des positions masquées où la meilleure prédiction
        de la tête MLM, $\arg\max_w \ell_{t,w}$, est le mot original.
  \item \textbf{Cosinus des paires similaires et dissimilaires}~: soit $\mathcal{S}$ les 338 paires de STS-B notées
        4 ou plus (similaires) et $\mathcal{D}$ les 308 paires notées 1 ou moins (dissimilaires). Pour une paire de
        phrases de plongements $(z, z')$, le notebook rapporte les cosinus moyens et leur écart,
        \[
          \overline{\cos}_{\mathcal{S}} = \frac{1}{|\mathcal{S}|} \sum_{(z, z') \in \mathcal{S}} \cos(z, z'), \qquad
          \overline{\cos}_{\mathcal{D}} = \frac{1}{|\mathcal{D}|} \sum_{(z, z') \in \mathcal{D}} \cos(z, z'), \qquad
          \text{écart} = \overline{\cos}_{\mathcal{S}} - \overline{\cos}_{\mathcal{D}}.
        \]
        Un bon plongement de phrase donne aux paires similaires un cosinus beaucoup plus élevé qu'aux paires
        dissimilaires~: un grand écart.
\end{itemize}

\begin{enumerate}
  \item \textbf{Entraînement (2 pts).} Incluez les courbes de perte (perte d'entraînement et de validation de chaque
  modèle en fonction des étapes d'entraînement).
  \begin{enumerate}
    \item L'un des trois modèles a-t-il surajusté~? Justifiez votre réponse à l'aide des courbes.
  \end{enumerate}

  \item \textbf{Niveau des mots et niveau des phrases (7 pts).} Incluez le tableau des résultats, le graphique de
  l'exactitude sur les mots masqués, le graphique des cosinus, et la sortie du texte à trous pour les deux phrases
  données et une phrase de votre choix.
  \begin{enumerate}
    \item Quel modèle a la meilleure exactitude sur les mots masqués~? Expliquez. [2 phrases max.]
    \item Lequel a le plus grand écart de cosinus~? Expliquez. [2 phrases max.]
  \end{enumerate}

  \item \textbf{Classification et apprentissage contrastif (3 pts).} Incluez la sortie de la recherche pour les deux
  requêtes données et une requête de votre choix.
  \begin{enumerate}
    \item Les deux modèles NLI ont été entraînés sur les mêmes phrases avec les mêmes étiquettes. Expliquez les
          différences d'écart de cosinus et de résultats de recherche. [2 phrases max.]
  \end{enumerate}

  \item \textbf{Pré-entraînement et affinage (2 pts).} En pratique, les encodeurs de phrases sont d'abord
  pré-entraînés avec MLM, puis affinés avec une perte contrastive. À partir de vos résultats et des données dont
  chaque objectif a besoin, expliquez ce que chaque étape apporte et pourquoi elles sont appliquées dans cet ordre.
  [2 phrases max.]
\end{enumerate}

\end{exercise}

\end{document}
